Decarbonization of energy-using sectors is essential for tackling climate change.
We use an ensemble of global integrated assessment models to assess CO2 emissions reduction potentials in buildings and transport, accounting for system interactions.
We focus on three intervention strategies with distinct emphases: reducing or changing activity, improving technological efficiency and electrifying energy end use. We find that these strategies can reduce emissions by 51–85\% in buildings and 37–91\% in transport by 2050 relative to a current policies scenario (ranges indicate model variability).
Electrification has the largest potential for direct emissions reductions in both sectors.
Interactions between the policies and measures that comprise the three strategies have a modest overall effect on mitigation potentials.
However, combining different strategies is strongly beneficial from an energy system perspective as lower electricity demand reduces the need for costly supply-side investments and infrastructure.
